% TeX root=../main.tex

\lipsum[1]

\clearpage{}
\textbf{Transliteração cirílica}

\begin{verse}
	\large
	\cirilex{i2 jutro probEzgu dzElo. vUstavU izide i@sI}\\
	\cirilex{i ide vU pusto mEsto. i2 tu molitvõ dEaSe.}\\
	\cirilex{i2 gUnaSẽ i simonU i iZe bEaxõ sU nimI.}\\
	\cirilex{i2 obrEtUSe i g@laSẽ emu. Eko vIsi iStõtU tebe.}\\
	\cirilex{i2 g@la imU idEmU vU bliZinẽjẽ vIsi i grady.}\\
	\cirilex{da i tu propovEmI.}\\
	\cirilex{na se bo izidU.}\\
	\cirilex{i2 bE propovEdajẽ na sUnUmiStixU ixU.}\\
	\cirilex{vU vUsei Galilei.}\\
	\cirilex{i2 bEsy izgonẽ.}\\
	\cirilex{i2 pride kU nemu prokaZenU molẽ i.}\\
	\cirilex{i2 na kolEnu padajẽ i g@lẽ emu.}\\
	\cirilex{Eko aSte xoSteSi moZeSi mẽ iStistiti.}\\
	\cirilex{i2@s Ze milosr@dovavU prosterU rõkõ kosnõ i.}\\
	\cirilex{i2 g@la emu xoStõ iSttisti sẽ.}\\
	\cirilex{i2 rekUS'u emu. abie otide prokaza otU nego. i2 CistU bystU.}\\
	\cirilex{i2 zaprEStI emu abie izgUna i.}\\
	\cirilex{i2 g@la emu bl'udi sẽ nikomuZe niCesoZe ne rIci.}\\
	\cirilex{nU SedU pokaZi sẽ arxiereovi.}\\
	\cirilex{i prinesi za oCiStenie tvoe.}\\
	\cirilex{eZe provEle mosi vU sUvEdEnie imU.}\\
	\cirilex{onU Ze iSedU naCẽtU propovEdati mUnogo.}\\
	\cirilex{i pronositi slovo.}\\
	\cirilex{Eko k tomu ne moZaaSe EvE vU gradU vIniti.}\\
	\cirilex{nU vInE vU pustExU mEstExU bE.}\\
	\cirilex{i2 prixoZdaxõ kU nemu otU vUsõdE.}
\end{verse}


\clearpage{}
\textbf{Transliteração latina}

\begin{verse}
	\large
	\ocstrans{i jutro probEzgu dzElo. vUstavU izide *i(su)sI}\\
	\ocstrans{i ide vU pusto mEsto. i tu molitvõ dEaSe.}\\
	\ocstrans{i gUnaSẽ i *simonU i iZe bEaxõ sU nimI.}\\
	\ocstrans{i obrEtUSe i g(lago)laSẽ emu. Eko vIsi iStõtU tebe.}\\
	\ocstrans{i g(lago)l imU idEmU vU bliZinẽjẽ vIsi i grady.}\\
	\ocstrans{da i tu propovEmI.}\\
	\ocstrans{na se bo izidU.}\\
	\ocstrans{i bE propovEdajẽ na sUnUmiStixU ixU.}\\
	\ocstrans{vU vUsei Galilei.}\\
	\ocstrans{i bEsy izgonẽ.}\\
	\ocstrans{i pride kU nemu prokaZenU molẽ i.}\\
	\ocstrans{i na kolEnu padajẽ i g(lago)lẽ emu.}\\
	\ocstrans{Eko aSte xoSteSi moZeSi mẽ iStistiti.}\\
	\ocstrans{*i(su)s Ze milosr(I)dovavU prosterU rõkõ kosnõ i.}\\
	\ocstrans{i g(lago)la emu xoStõ iSttisti sẽ.}\\
	\ocstrans{i rekUS'u emu. abie otide prokaza otU nego. i CistU bystU.}\\
	\ocstrans{i zaprEStI emu abie izgUna i.}\\
	\ocstrans{i g(lago)la emu bl'udi sẽ nikomuZe niCesoZe ne rIci.}\\
	\ocstrans{nU SedU pokaZi sẽ arxiereovi.}\\
	\ocstrans{i prinesi za oCiStenie tvoe.}\\
	\ocstrans{eZe provEle Mosi vU sUvEdEnie imU.}\\
	\ocstrans{onU Ze iSedU naCẽtU propovEdati mUnogo.}\\
	\ocstrans{i pronositi slovo.}\\
	\ocstrans{Eko k tomu ne moZaaSe EvE vU gradU vIniti.}\\
	\ocstrans{nU vInE vU pustExU mEstExU bE.}\\
	\ocstrans{i prixoZdaxõ kU nemu otU vUsõdE.}
\end{verse}


\clearpage{}
\textbf{Tradução}

\begin{verse}
	\large
	De madrugada,estando ainda escuro, Ele se levanta\\
	e retirou-se para um lugar deserto e ali orava,\\
	Simão e os seus companheiros o procuravam ansiosos\\
	e, quando o acharam, disseram-lhe: “Todos te procuram”.\\
	Disse-lhes: ``Vamos a outros lugares, às aldeias da vizinhança,\\
	a fim de pregar também ali,\\
	pois foi para isso que Eu sai''.\\
	E foi pregando em  sinagogas\\
	por toda a Galiléia,\\
	e expulsando os demônios.\\
	Um leproso foi até Ele, implorando-lhe,\\
	de joelhos dizendo-lhe:\\
	``Se queres, tens o poder de purificar-me''.\\
	Movido de compaixão, estendeu a mão, toucou-o\\
	e disse-lhe: ``Eu quero, sê purificado''.\\
	E logo a lepra o deixou e ficou purificado.\\
	Advertindo-o severamente, despediu-o logo.\\
	Dizendo-lhe: ``Não digas nada a ninguém;\\
	mas vai mostrar-te ao sacerdote\\
	e oferece por tua purificação\\
	o que Moisés prescreveu, para que lhes sirva de prova''.\\
	Ele, porém, assim que partiu, começou a proclamar ainda mais\\
	e a divulgar a notícia,\\
	de modo que Jesus já não podia entrar publicamente numa cidade:\\
	permanecia fora, em lugares desertos.\\
	E de toda parte vinham procurá-lo.
\end{verse}


\clearpage{}
\textbf{Original grego}

\begin{verse}
	\large
	Καὶ πρωῒ ἔννυχα λίαν ἀναστὰς ἐξῆλθεν\\
	καὶ ἀπῆλθεν εἰς ἔρημον τόπον, κἀκεῖ προσηύχετο.\\
	καὶ κατεδίωξεν αὐτὸν Σίμων καὶ οἱ μετ’ αὐτοῦ,\\
	καὶ εὗρον αὐτὸν καὶ λέγουσιν αὐτῷ ὅτι Πάντες ζητοῦσίν σε.\\
	καὶ λέγει αὐτοῖς· Ἄγωμεν ἀλλαχοῦ εἰς τὰς ἐχομένας κωμοπόλεις,\\
	ἵνα καὶ ἐκεῖ κηρύξω·\\
	εἰς τοῦτο γὰρ ἐξῆλθον.\\
	καὶ ἦλθεν κηρύσσων εἰς τὰς συναγωγὰς αὐτῶν\\
	εἰς ὅλην τὴν Γαλιλαίαν\\
	καὶ τὰ δαιμόνια ἐκβάλλων.\\
	Καὶ ἔρχεται πρὸς αὐτὸν λεπρὸς παρακαλῶν αὐτὸν\\
	καὶ γονυπετῶν λέγων αὐτῷ\\
	ὅτι Ἐὰν θέλῃς δύνασαί με καθαρίσαι.\\
	καὶ σπλαγχνισθεὶς ἐκτείνας τὴν χεῖρα αὐτοῦ ἥψατο\\
	καὶ λέγει αὐτῷ Θέλω, καθαρίσθητι.\\
	καὶ εὐθὺς ἀπῆλθεν ἀπ’ αὐτοῦ ἡ λέπρα, καὶ ἐκαθερίσθη.\\
	καὶ ἐμβριμησάμενος αὐτῷ εὐθὺς ἐξέβαλεν αὐτόν,\\
	καὶ λέγει αὐτῷ Ὅρα μηδενὶ μηδὲν εἴπῃς,\\
	ἀλλὰ ὕπαγε σεαυτὸν δεῖξον τῷ ἱερεῖ\\
	καὶ προσένεγκε περὶ τοῦ καθαρισμοῦ σου\\
	ἃ προσέταξεν Μωϋσῆς εἰς μαρτύριον αὐτοῖς.\\
	ὁ δὲ ἐξελθὼν ἤρξατο κηρύσσειν πολλὰ\\
	καὶ διαφημίζειν τὸν λόγον,\\
	ὥστε μηκέτι αὐτὸν δύνασθαι φανερῶς εἰς πόλιν εἰσελθεῖν,\\
	ἀλλ’ ἔξω ἐπ’ ἐρήμοις τόποις ἦν·\\
	καὶ ἤρχοντο πρὸς αὐτὸν πάντοθεν.
\end{verse}
